\pagestyle{fancy}
\fancyhf{}
\renewcommand{\headrulewidth}{1pt}
\lhead{\color{black}{ЮТКБ.656122.611-01.01 РЭ}}
\chead{\color{black}{Редакция \rev~от~\revdate}}
\rhead{\color{unigreen}\colondevicetype}
\lfoot{\color{unigreen}{Руководство по эксплуатации }}
\rfoot{{\thepage}}
\renewcommand{\footrulewidth}{1pt}

%\vskip 10mm

\pagecolor{white}

\noindent
$\copyright$ \the\year{} Юнител Инжиниринг
\hfill
\parbox{0.3\textwidth}{\raggedleft Редакция: \textcolor{unigray}{\rev}}

\vskip 3mm
\noindent
Москва
\hfill
\parbox{0.3\textwidth}{\raggedleft Дата: \textcolor{unigray}{\revdate}}
\vskip 5mm

\par
Настоящее руководство по эксплуатации (РЭ) предназначено для ознакомления с основными параметрами, принципом действия, конструкцией, правилами эксплуатации и обслуживания микропроцессорного устройства серии «ЮНИТ-М500-ВЧ», именуемого в дальнейшем «Устройство».

\par
Устройство соответствует:
\begin{itemize}
\item ТУ 27.12.23-609.61775353-2024 Устройство защиты и автоматики серии ЮНИТ-М500;
\item требованиям ТР ТС «О безопасности низковольтного оборудования» ТР ТС 004/2011 при выполнении требований: ГОСТ 12.2.007.0-75 «Система стандартов безопасности труда. Изделия электротехнические. Общие требования безопасности», ГОСТ 12.1.004-91 «Система стандартов безопасности труда. Пожарная безопасность. Общие требования», ГОСТ 14254-96 «Степени защиты, обеспечиваемые оболочками (код IP)»; 
\item требованиям ТР ТС «Электромагнитная совместимость технических средств» ТР ТС 020/2011 при выполнении требований ГОСТ Р 51317.6.5-2006 (МЭК 61000-6-5:2001) «Совместимость технических средств электромагнитная. Устойчивость к электромагнитным помехам технических средств, применяемых на электростанциях и подстанциях. Требования и методы испытаний»;
\item РД 34.35.310-97 (СО 34.35.310-97) «Общие технические требования к микропроцессорным устройствам защиты и автоматики энергосистем»;
\item Приказ министерства энергетики Российской Федерации от 13 февраля 2019 года № 101 (с изменениями на 10 июля 2020 года);
\item СТО 56947007-29.120.70.241-2017 с изменениями от 11.12.2019 «Технические требования к микропроцессорным устройствам РЗА. Стандарт организации ПАО ''ФСК ЕЭС'' 2017 год»;
\item СТО 34.01-6.2-002-2024 «Контроллеры присоединения, преобразователи дискретных сигналов. Типовые технические требования»;
\item CTO 34.01-21-004-2019 «Цифровой питающий центр. Требования к технологическому проектированию цифровых подстанций напряжением 110-220 кВ и узловых цифровых подстанций напряжением 35 кВ»;
\item СТО 56947007-29.240.10.256-2018 «Технические требования к аппаратно-программным средствам и электротехническому оборудованию ЦПС»;
\item CTO 34.01-21-005-2019 «Цифровая электрическая сеть. Требования к проектированию цифровых распределительных электрических сетей 0,4-220 кВ»;
\item \sloppy ОТТ-29.020.00-КТН-009-15 «Магистральный трубопроводный транспорт нефти и нефтепродуктов. Микропроцессорные устройства РЗА подстанций 35-220 кВ и распределительных устройств 6(10) кВ. Общие технические требования»;
\item CTO Газпром 2-1.11-661-2012 «Цифровые устройства релейной защиты и автоматики для систем электроснабжения. Технические требования»;
\item СТО 56947007-25.040.30.309-2020 с изменениями от 24.03.2023 «Корпоративный профиль МЭК 61850 ПАО ''ФСК ЕЭС''»;
\item СТО 56947007-33.040.20.282-2019 Типовые шкафы ШЭТ РЗА ЛЭП 110 – 750 кВ. Архитектура I типа;
\item СТО 56947007-33.040.20.283-2019 Типовые шкафы ШЭТ РЗА ЛЭП 110 – 750 кВ. Архитектура II типа.
\end{itemize}

%\medskip
\par
К эксплуатации устройства допускаются лица, изучившие настоящее РЭ и прошедшие проверку знаний правил техники безопасности и эксплуатации электроустановок электрических станций и подстанций.
%\par
%Для проведения работ с микропроцессорным устройством допускаются лица, изучившие настоящее руководство, прошедшие проверку знаний правил техники безопасности и эксплуатации электроустановок, электрических станций и подстанций, и имеющие группу по электробезопасности не ниже III по работе с электроустановками напряжением до 1000~В.

%\medskip

\par
Необходимые параметры и надежность работы устройств в течение срока службы обеспечиваются качеством разработки и изготовления, соблюдением условий транспортирования, хранения, монтажа, наладки и обслуживания, поэтому выполнение всех требований настоящего РЭ является обязательным.


%\medskip
\par
Компания Юнител Инжиниринг оставляет за собой авторские права на данный документ и на информацию, содержащуюся в нем, включая права на использование патентов. Копирование, использование и передача информации третьим лицам без письменного разрешения компании категорически запрещены.

%\medskip
\par
По всем интересующим вопросам можно обращаться:
\begin{itemize}
\item по адресу электронной почты: \href{mailto:rza@uni-eng.ru}{\color{unigreen}\uline{rza@uni-eng.ru}};
\item по телефону: +7 (495) 165-99-98 (доб. 2561).
\end{itemize}
Также с нашей продукцией можно ознакомиться на сайте: \href{https://uni-eng.ru}{\color{unigreen}\uline{https://uni-eng.ru}}.